\section{Extended Comparative Statics and Counterexamples}
\label{sec:or-extended-comparative-statics}

These results are supporting interpretations of the finite dynamic model.
They are not required by the tagged-cover reduction and are not used to select
financial retention decisions.

\input{sections/09_comparative_statics_counterexamples}
\subsection{Auxiliary finite-model illustration}

The following exact construction isolates the distinction between operating
profiles and generative modules in the main paper's finite model.  It is an
illustration rather than an additional assumption or theorem.

\begin{example}[A two-module generator]\label{ex:modular-generator}
Let \(X=\{\mathrm{down},\mathrm{up}\}\), let
\(\mu_{b^-}=(3/4,1/4)\) and \(\mu_{b^+}=(1/4,3/4)\), and take identity
closure.  Strategies \(s_T\) and \(s_H\) carry modules \(m_T\) and \(m_H\).
A project \(q\) requires both.  When
\(\{m_T,m_H\}\subseteq C\), set
\[
  G(q,b,C)(\operatorname{some}(c))=\frac23,\qquad
  G(q,b,C)(\operatorname{none})=\frac13,\qquad
  \nu(q,b,C,c)=\frac34;
\]
otherwise set \(G(q,b,C)=\delta_{\operatorname{none}}\).  Retaining both
modules therefore admits \(c\) with probability \(1/2\), while deleting the
sole carrier of either module makes admission impossible.  If
\(u_c(\mathrm{down})=-1\) and \(u_c(\mathrm{up})=2\), then
\(j_c(b^-)=-1/4\) and \(j_c(b^+)=5/4\): current operating payoff depends on
the filtered state, while candidate availability depends on retained closure.
For a positive-duration instance, take \(d_q=2\), initiation cost \(1/10\),
and operation flag \(o_q=1\).
\end{example}

\subsection{Conditional replacement under a capacity bound}
\label{app:conditional-replacement}

The supporting question is what an incumbent should release when an
outer-certified candidate \(c\notin L\) becomes available for retention and
capacity binds.  This formulation is conditional on outer certification and
is separate from the raw generation, verification, and admission law in the
main paper.

Let the incumbent satisfy \(W(L)\le B\), let \(w_c\le B\), and permit deletion
of active incumbents \(D\subseteq L\setminus\{s_0\}\).  Additivity makes the
required release
\[
 r_c(B):=[W(L)+w_c-B]_+,
\]
so the feasible deletion family is
\[
 \mathcal D_c(B)
 :=\{D\subseteq L\setminus\{s_0\}:W(D)\ge r_c(B)\}.
\]
The conditional replacement problem is
\begin{equation}\label{eq:optional-replacement-problem}
 \operatorname*{maximize}_{D\in\mathcal D_c(B)}
 V_\theta\bigl(b,(L\setminus D)\cup\{c\}\bigr).
\end{equation}

For interpretation, define the unconstrained candidate gain and displacement
loss
\begin{align*}
 G_c(b,L)&:=V_\theta(b,L\cup\{c\})-V_\theta(b,L),\\
 \ell_c(D)&:=V_\theta(b,L\cup\{c\})
 -V_\theta\bigl(b,(L\setminus D)\cup\{c\}\bigr).
\end{align*}
Then the best net admission value is the definition-level identity
\[
 A_c(b,L,B)=G_c(b,L)-\min_{D\in\mathcal D_c(B)}\ell_c(D).
\]
Under common insertion and the productive factorization, a deletion safe
before admission is sufficient for zero displacement loss, but it is not
necessary: the candidate may replace a capability carried by a deleted
incumbent.  This supporting formulation is not part of the Lean-verified
theorem set, and it makes no constrained--penalized equivalence claim.

\subsection{Secondary sensitivity results}
\label{or:aor-secondary-comparative-statics}

This section retains sensitivity calculations that are useful for interpreting
the economic model but are not needed for the safe-compression reduction.
Each calculation holds along the named parameter path or conditional on the
displayed initial belief and fixed action. Unnamed primitives are fixed.
Finiteness is used for termwise differentiation and finite action
maximization. No closure-detectability or stationary-contraction argument is
used.

\subsection{Bridge-margin sensitivity}

Write
\begin{equation}
 A_{\mathrm{br}}:=\beta^d\rho^d\pi C,
 \qquad M_{\mathrm{br}}:=A_{\mathrm{br}}-\kappa,
 \qquad m_{\mathrm{br}}:=\frac{M_{\mathrm{br}}}{A_{\mathrm{br}}}.
 \label{eq:or-bridge-margin}
\end{equation}
On the declared real extension with fixed integer \(d\), positive gross
coordinates, and \(M_{\mathrm{br}}>0\), realized bridge value equals the
margin. Its named-coordinate elasticities are
\begin{equation}
 (\varepsilon_\beta^M,\varepsilon_\rho^M,
   \varepsilon_\pi^M,\varepsilon_C^M,\varepsilon_\kappa^M)
 =\left(\frac d{m_{\mathrm{br}}},\frac d{m_{\mathrm{br}}},
   \frac1{m_{\mathrm{br}}},\frac1{m_{\mathrm{br}}},
   -\frac{1-m_{\mathrm{br}}}{m_{\mathrm{br}}}\right).
 \label{eq:or-bridge-elasticities}
\end{equation}
Threshold fragility is governed by the normalized positive margin
\(m_{\mathrm{br}}\). The magnitudes diverge as
\(m_{\mathrm{br}}\downarrow0\), but not merely because
\(M_{\mathrm{br}}\downarrow0\): if cost is zero and gross scale vanishes,
then \(m_{\mathrm{br}}=1\). At \(M_{\mathrm{br}}=0\), bridge-value
elasticity is undefined. At and below that threshold, the signed margin,
action region, or exact finite changes in realized value replace a percentage
elasticity.

For the derivation, condition on the canonical project and hold every unnamed
coordinate fixed. Differentiating on the positive real extension gives
\[
 \beta\,\partial_\beta M_{\mathrm{br}}=dA_{\mathrm{br}},
 \quad
 \rho\,\partial_\rho M_{\mathrm{br}}=dA_{\mathrm{br}},
 \quad
 \pi\,\partial_\pi M_{\mathrm{br}}
 =C\,\partial_C M_{\mathrm{br}}=A_{\mathrm{br}},
 \quad
 \kappa\,\partial_\kappa M_{\mathrm{br}}=-\kappa.
\]
On \(M_{\mathrm{br}}>0\), realized value
\([M_{\mathrm{br}}]_+\) equals the signed margin. Dividing by
\(M_{\mathrm{br}}\) and using
\(m_{\mathrm{br}}=M_{\mathrm{br}}/A_{\mathrm{br}}\) gives
\eqref{eq:or-bridge-elasticities}. At the threshold, the positive-part value
has no ordinary percentage elasticity. The costless path
\(\kappa=0\), \(A_{\mathrm{br}}\downarrow0\) keeps
\(m_{\mathrm{br}}=1\), so a vanishing net level alone does not imply the
threshold-fragility limit.

\subsection{Operational and generative elasticity contributions}

Let one named positive scalar \(x\) vary the same primitive in total insertion
value \(I(x)\), operational value \(O(x)\), and generative value \(G(x)\),
holding every unlisted primitive fixed and remaining on a differentiable
branch. Differentiating \(I=O+G\) gives
\[
 xI'(x)=xO'(x)+xG'(x).
\]

\begin{proposition}[Operational--generative elasticity contributions]
\label{prop:or-channel-elasticity}
At \(x_0>0\), suppose \(I(x)=O(x)+G(x)\) throughout a neighborhood of
\(x_0\), and the three channel paths are differentiable there along the same
named perturbation. If \(I(x_0)>0\), define
\[
 C_x^{\mathrm{op}}:=\frac{x_0O'(x_0)}{I(x_0)},
 \qquad
 C_x^{\mathrm{gen}}:=\frac{x_0G'(x_0)}{I(x_0)}.
\]
Then
\[
 \varepsilon_x^I=C_x^{\mathrm{op}}+C_x^{\mathrm{gen}}.
\]
Only when both channel levels are positive may this be written as
\[
 \varepsilon_x^I
 =\frac{O}{I}\varepsilon_x^{\mathrm{op}}
  +\frac{G}{I}\varepsilon_x^{\mathrm{gen}}.
\]
\end{proposition}

The contribution form remains meaningful when one channel level is zero; the
corresponding ordinary component log elasticity does not. If total value is
zero, no total elasticity or total-normalized contribution is reported. The
scaled derivative identity above, or the exact finite-change identity
\(\Delta I=\Delta O+\Delta G\), remains well defined.

Along one fixed differentiable scalar path, differentiating the level identity
and applying common scaling gives
\[
 xI'(x)=xO'(x)+xG'(x).
\]
If \(I(x)>0\), division by the total level gives the two signed contributions
in \cref{prop:or-channel-elasticity}. Rewriting a contribution as
\[
 \frac{xO'}I=\frac OI\frac{xO'}O
\]
requires \(O>0\), and the analogous rewrite requires \(G>0\). If a component
vanishes, its scaled derivative remains available but its log elasticity does
not. If \(I=0\), neither total elasticity nor a total-normalized contribution
is defined; the level derivative or exact finite-change identity is used.

\begin{table}[ht]
\centering
\scriptsize
\setlength{\tabcolsep}{2.5pt}
% Generated by julia/scripts/run_unified_resource_benchmark.jl
\begin{tabular}{llrrrrrr}
\toprule
Belief & Library & $V$ & $V^{\mathrm{op}}$ & $V^{\mathrm{gen}}$ & $D_\beta$ & $E^{\mathrm{op}}_\beta$ & $E^{\mathrm{gen}}_\beta$ \\
\midrule
low & I & $\frac{115}{288}$ & $0$ & $\frac{115}{288}$ & $\frac{10777}{2070}$ & $--$ & $\frac{10777}{2070}$ \\
low & I+A & $1$ & $0$ & $1$ & $\frac{56}{15}$ & $--$ & $\frac{56}{15}$ \\
low & I+D & $\frac{14}{3}$ & $\frac{14}{3}$ & $0$ & $\frac{25}{21}$ & $\frac{25}{21}$ & $--$ \\
high & I & $\frac{23}{288}$ & $0$ & $\frac{23}{288}$ & $\frac{14089}{2070}$ & $--$ & $\frac{14089}{2070}$ \\
high & I+A & $\frac{7}{6}$ & $0$ & $\frac{7}{6}$ & $\frac{53}{15}$ & $--$ & $\frac{53}{15}$ \\
high & I+D & $\frac{22}{3}$ & $\frac{22}{3}$ & $0$ & $\frac{29}{33}$ & $\frac{29}{33}$ & $--$ \\
\bottomrule
\end{tabular}

\caption{Exact productive value, channel levels, and fixed-policy discount
elasticity in the canonical benchmark. A dash marks an undefined ordinary
component elasticity at a zero component level; it is never recorded as
zero. Scaled contributions are evaluated directly.}
\label{tab:or-canonical-channel-elasticities}
\end{table}

\subsection{Fixed-exposure innovation duration}

For a fixed nonnegative, nonzero exposure sequence
\(z_0,\ldots,z_{H-1}\) and effective discount \(\alpha>0\), let
\begin{equation}
 \Psi_H(\alpha):=\sum_{t=0}^{H-1}\alpha^tz_t,
 \quad \omega_t:=\frac{\alpha^tz_t}{\Psi_H(\alpha)},
 \quad D_\Psi:=\sum_t t\omega_t,
 \quad C_\Psi:=\sum_t\omega_t(t-D_\Psi)^2.
 \label{eq:or-innovation-potential-duration}
\end{equation}
Because the sequence is nonzero, \(\Psi_H(\alpha)>0\). The exact finite
scaled-derivative identities are
\begin{equation}
 \alpha\Psi_H'(\alpha)=D_\Psi\Psi_H(\alpha),
 \qquad
 \alpha D_\Psi'(\alpha)=C_\Psi\geq0.
 \label{eq:or-innovation-duration}
\end{equation}
Equivalently, \(D_\Psi\) is the derivative of \(\log\Psi_H\) with respect to
\(\log\alpha\), and \(C_\Psi\) is the corresponding derivative of
\(D_\Psi\). These are interpretations of the displayed scaled identities.
Innovation duration is the contribution-weighted date at which fixed
uncovered value arrives, and innovation convexity is its timing variance. A
higher effective discount shifts weight toward later exposures and weakly
raises duration. This is not primitive project delay, time to admission, or a
convexity theorem for an optimized Bellman value whose exposure sequence or
selected action changes. If exposure is zero, duration and log elasticity are
undefined; level contributions or finite differences are used instead.

For the derivation, introduce
\[
 N_1(\alpha):=\sum_t t\alpha^tz_t,
 \qquad
 N_2(\alpha):=\sum_t t^2\alpha^tz_t.
\]
Termwise finite differentiation gives
\[
 \alpha\Psi_H'(\alpha)=N_1(\alpha),
 \qquad
 \alpha N_1'(\alpha)=N_2(\alpha).
\]
Since \(\Psi_H>0\), \(D_\Psi=N_1/\Psi_H\), and the quotient rule yields
\[
 \alpha D_\Psi'
 =\frac{N_2\Psi_H-N_1^2}{\Psi_H^2}
 =\sum_t\omega_t(t-D_\Psi)^2=C_\Psi\geq0.
\]
The nonnegative curvature is curvature of \(\log\Psi_H\) in
\(\log\alpha\) for a fixed exposure sequence. Changing project duration,
belief dynamics, completion laws, or the selected action changes the sequence
and lies outside this derivative.

\subsection{Finite switch and one-sided-change conventions}

Within a strictly dominant price branch, the penalized envelope inherits the
branch slope \(-W(L')\). A pairwise crossing is only a candidate switch; it
becomes a resource-price breakpoint only when both branches are globally
active. At an unequal-burden tie the optimizer is set-valued, so one-sided
slopes or cross-breakpoint finite changes replace a derivative selected by an
unstated tie breaker.

For a finite library family, record the best-versus-second-best objective
margin; for a Bellman node, record the corresponding action-value margin. A
positive margin measures distance from a tie in value space. A zero margin
flags a possible switch, but distinct libraries or actions may generate
identical branches, so it is not sufficient for a change of optimizer. At an
actual switch with positive value, one-sided or arc elasticities may be
reported; at zero value, only level slopes and finite differences are defined.


\subsection{Finite capacity and resource-price geometry}
\label{sec:or-capacity-price-geometry}

Let \(\mathfrak L_\theta\) be a nonempty finite family of outer-certified
libraries, let \(W(L')\geq0\) be retention burden, and let
\(V_\theta(b,L')\) be productive value under the unchanged inner process.
Define
\[
 V_\theta^\star(b;B)
 :=\max_{L'\in\mathfrak L_\theta:\,W(L')\leq B}V_\theta(b,L').
\]

\begin{theorem}[Finite hard-capacity retention]
\label{thm:or-capacity-value}
Suppose \(\mathfrak L_\theta\) is finite and contains the zero-burden
inactive-only library. For every \(B\geq0\), the maximum is attained. The
function \(B\mapsto V_\theta^\star(b;B)\) is nondecreasing, is constant
between consecutive attainable burdens, and can change only at an attainable
burden. Every exact forward capacity increment is nonnegative. Neither
uniqueness nor inclusion nesting of maximizing libraries follows.
\end{theorem}

\begin{proof}
For fixed \(B\), the feasible family is a nonempty subset of a finite set, so
the maximum is attained. If capacity rises, feasible families are nested,
which proves monotonicity. Between consecutive values in the finite set
\(\{W(L'):L'\in\mathfrak L_\theta\}\), the feasible family does not change.
The hypotheses impose neither strict value comparisons nor inclusion relations
among maximizers.
\end{proof}

Discreteness matters. With two unit-burden capabilities that produce no value
separately and value one jointly, successive attainable values are
\((0,0,1)\). The capacity function is not generally concave, marginal capacity
increments need not diminish, and a finite nonconcave capacity optimum need
not be supported by a scalar price.

For \(\lambda\geq0\), define the penalized problem
\[
 J_\theta^\star(b;\lambda)
 :=\max_{L'\in\mathfrak L_\theta}
 \{V_\theta(b,L')-\lambda W(L')\}.
\]

\begin{theorem}[Finite resource-price envelope]
\label{thm:or-penalized-envelope}
For a fixed nonempty finite eligible family, the maximum is attained. The
optimized value is continuous, nonincreasing, convex, and piecewise affine in
the real price \(\lambda\geq0\), with switching candidates among finitely many
pairwise branch intersections. If \(0\leq\lambda_1<\lambda_2\), every optimizer
at \(\lambda_2\) has weakly lower burden than every optimizer at
\(\lambda_1\). Wherever one library is strictly dominant, the local slope is
\(-W(L')\).
\end{theorem}

\begin{proof}
Each eligible library supplies the affine branch
\(V_\theta(b,L')-\lambda W(L')\). The maximum of finitely many continuous
affine functions is continuous, convex, and piecewise affine; nonnegative
burdens make every branch nonincreasing. Let \(L_i\) optimize at
\(\lambda_i\), \(i=1,2\). Adding the two optimality inequalities and using
\(\lambda_2>\lambda_1\) gives \(W(L_2)\leq W(L_1)\). On an interval with a
unique active branch, differentiation gives slope \(-W(L')\).
\end{proof}

For two unequal-burden libraries, their value difference divided by their
burden difference is only a candidate crossing price; it is an actual
breakpoint only if both branches reach the upper envelope. At a tie the
optimizer is set-valued. The burden ordering across prices does not imply that
raw strategy identities are nested.
